\documentclass[11pt]{article}

\usepackage[margin=1in]{geometry}
\usepackage[T1]{fontenc}
\usepackage{lmodern}
\usepackage{amsmath,amssymb}
\usepackage{booktabs}
\usepackage{tabularx}
\usepackage{enumitem}
\usepackage[hidelinks]{hyperref}
\usepackage{xurl}

\newcommand{\code}[1]{\texttt{#1}}
\newcommand{\pred}[1]{\texttt{#1}}

\title{Methods: Predicate frequencies in SemMedDB-KGX and\\
recall against CTD curated chemical--disease associations}
\author{Reuben Thomas}
\date{\today}

\begin{document}
\maketitle

\section{Overview}

We performed two analyses on the Biolink-normalized Knowledge Graph Exchange
(KGX) release of the Semantic MEDLINE Database (SemMedDB)~\cite{kilicoglu2012,semmedkgx}:
(i)~the frequency of each (subject category, predicate, object category)
triple across the edge set, and (ii)~the recall of SemMedDB chemical--disease
edges relative to the manually curated chemical--disease associations in the
Comparative Toxicogenomics Database (CTD)~\cite{ctd}, evaluated only on
PubMed articles indexed by both resources. Both analyses were implemented as
SQL queries in DuckDB~\cite{duckdb} and orchestrated with
Snakemake~\cite{snakemake}.

\section{Data sources}

\paragraph{SemMedDB-KGX.}
We used the ``uncapped'' SemMedDB-KGX archive
\code{kg2.10.3\_semmeddb\_dogpark\_uncapped\_2026\_04\_07.tar.gz}
from the Translator-CATRAX SemMedDB-KGX project~\cite{semmedkgx}, and extracted two
files from it:
\begin{itemize}[nosep]
  \item \code{normalized\_nodes.jsonl}: one record per node, giving its CURIE
        (a UMLS CUI), an ordered list of Biolink categories~\cite{biolink},
        and the node normalizer's \code{equivalent\_identifiers}
        cross-references;
  \item \code{normalized\_edges.jsonl}: one record per
        (subject, predicate, object) edge, with Biolink predicates and the list
        of supporting PubMed identifiers (\code{publications}).
\end{itemize}
This edge set is a filtered subset of raw SemMedDB. Upstream it was filtered
by minimum supporting-publication count, by Biolink domain/range constraints,
and by removing predicates excluded by BioThings Explorer (see the
SemMedDB-KGX documentation). So all frequencies and recall estimates below
describe this curated edge set, not the raw SemMedDB \code{PREDICATION}
table.

\paragraph{CTD.}
We downloaded CTD's curated chemical--disease associations file
(\code{CTD\_curated\_chemicals\_diseases.csv.gz}) from
\url{https://ctdbase.org/reports/}. Each row gives a chemical (a MeSH
accession), a disease (a MeSH or OMIM CURIE), a \code{DirectEvidence} value
(\code{therapeutic} or \code{marker/mechanism}), and a pipe-delimited list of
the PubMed IDs that support that evidence type. CTD does not version this file
at the URL, so the release we used is the one current on the download date.

\section{Category assignment}

Each node was assigned its \emph{most specific} Biolink category: the first
element of its \code{category} list, following the KGX convention of ordering
categories from most to least specific. The \code{biolink:} prefix was removed
from categories and predicates.

\section{Predicate frequency estimation}

We joined every edge to the most specific category of its subject and object
nodes, and grouped the edges by
(subject category, predicate, object category). Each group reports two
counts:
\begin{itemize}[nosep]
  \item $n_{\text{edges}}$: the number of distinct edges in the group;
  \item $n_{\text{publications}}$: the sum, over those edges, of the length of
        each edge's \code{publications} list. This is a count of edge--article
        support instances, not of distinct articles.
\end{itemize}
Edges whose subject or object had no node record are excluded by the inner
join. We also reported counts separately for a configured set of triples of
interest (for example, \code{Procedure}--\pred{has\_input}--\code{PhysicalEntity}
and \code{Disease}--\pred{located\_in}--\code{GrossAnatomicalStructure}). A
triple that does not occur is reported with zero counts.

\section{Recall against CTD}

\subsection{Mapping SemMedDB concepts to MeSH}

SemMedDB identifies concepts by UMLS CUI, while CTD uses MeSH. We mapped CUIs
to MeSH using the \code{MESH:} entries in each node's
\code{equivalent\_identifiers}. Because this mapping comes from the node
normalizer's existing cross-references, it does not require a UMLS
Metathesaurus (MRCONSO) license. A CUI with several MeSH cross-references
keeps all of them, so the mapping can be many-to-many.

Chemical-like and disease-like nodes were defined by their most specific
category:
\begin{itemize}[nosep]
  \item chemical: \code{SmallMolecule}, \code{ChemicalEntity}, \code{Drug},
        \code{ChemicalMixture};
  \item disease: \code{Disease}, \code{PhenotypicFeature},
        \code{DiseaseOrPhenotypicFeature}.
\end{itemize}

\subsection{SemMedDB chemical--disease assertions}

We kept only edges with a chemical-like subject and a disease-like object, so
edges asserted in the opposite direction (disease $\rightarrow$ chemical) are
not considered. Each retained edge was expanded over its \code{publications}
list, and its subject and object were mapped to MeSH. This gives the set of
distinct tuples
\[
  S = \{(c, p, d, m)\},
\]
where $c$ is a chemical MeSH ID, $p$ a Biolink predicate, $d$ a disease MeSH
ID and $m$ a PubMed ID. Subjects or objects without a MeSH cross-reference
drop out at this step.

\subsection{CTD reference assertions}

Each CTD row was expanded over its \code{DirectEvidence} values and its PubMed
IDs, and the chemical accession was prefixed with \code{MESH:}. This gives the
set of distinct tuples
\[
  C = \{(c, d, e, m)\},
\]
where $e \in \{\code{therapeutic}, \code{marker/mechanism}\}$. Rows whose
disease is an OMIM identifier (not MeSH) cannot match SemMedDB's MeSH-mapped
diseases, so we excluded them and recorded how many there were.

\subsection{Common-article anchor}

A CTD assertion can only be recovered by SemMedDB if SemMedDB processed the
same article, so we restricted the evaluation to articles both resources
cover. We defined the common PubMed set as
\[
  M = \pi_m(C) \cap \pi_m(S),
\]
where $\pi_m$ is the set of distinct PubMed IDs in a table, and $\pi_m(S)$ is
taken over chemical--disease edges only. The evaluation set is
$C_M = \{(c,d,e,m) \in C : m \in M\}$, and SemMedDB is restricted the same way
to $S_M$.

\subsection{Matching criteria}

Each CTD assertion $(c,d,e,m) \in C_M$ was scored against two criteria of
increasing strictness:
\begin{description}[nosep]
  \item[Entity match:] SemMedDB asserts \emph{some} relation between $c$ and
        $d$ in article $m$:
        $\exists\, p : (c,p,d,m) \in S_M$.
  \item[Predicate match:] SemMedDB asserts such a relation with a predicate
        consistent with CTD's evidence type:
        $\exists\, p \in P(e) : (c,p,d,m) \in S_M$.
\end{description}
The allowed predicate sets $P(e)$ are given in
Table~\ref{tab:predmap}. We chose them conservatively: generic predicates that
do not reflect curated evidence (\pred{related\_to}, \pred{coexists\_with},
\pred{located\_in}, \pred{produces}, \pred{has\_part}, \pred{derives\_from},
\pred{occurs\_in}, \pred{precedes}) count toward neither evidence type. For
each CTD assertion we also recorded the full list of SemMedDB predicates that
link the same chemical, disease and article, so that alternative predicate
mappings can be assessed afterwards.

\begin{table}[h]
  \centering
  \caption{SemMedDB predicates counted as consistent with each CTD
           \code{DirectEvidence} type.}
  \label{tab:predmap}
  \begin{tabularx}{0.9\linewidth}{@{}lX@{}}
    \toprule
    CTD \code{DirectEvidence} & Allowed SemMedDB predicates $P(e)$ \\
    \midrule
    \code{therapeutic}      & \pred{treats\_or\_applied\_or\_studied\_to\_treat} \\
    \code{marker/mechanism} & \pred{causes}, \pred{affects}, \pred{disrupts},
                              \pred{predisposes\_to\_condition},
                              \pred{exacerbates\_condition},
                              \pred{interacts\_with} \\
    \bottomrule
  \end{tabularx}
\end{table}

\subsection{Recall}

For each evidence type $e$, and for all types pooled, recall is the fraction
of CTD assertions in the evaluation set that are matched:
\[
  R^{\text{entity}}_e =
    \frac{\bigl|\{x \in C_M^{(e)} : x \text{ entity-matched}\}\bigr|}
         {\bigl|C_M^{(e)}\bigr|},
  \qquad
  R^{\text{predicate}}_e =
    \frac{\bigl|\{x \in C_M^{(e)} : x \text{ predicate-matched}\}\bigr|}
         {\bigl|C_M^{(e)}\bigr|},
\]
where $C_M^{(e)}$ is the subset of $C_M$ with evidence type $e$. The unit of
evaluation is a (chemical, disease, evidence type, article) tuple. In the
pooled estimate, a chemical--disease pair supported by the same article under
both evidence types therefore counts twice. By construction,
$R^{\text{predicate}}_e \le R^{\text{entity}}_e$.

\subsection{Diagnostics}

We also report: the number of common PubMed IDs and the number of distinct
PubMed IDs on each side; the number of CTD rows and distinct
chemical--disease pairs excluded for OMIM disease identifiers; and the
fraction of chemical-subject and disease-object CUIs (among chemical--disease
edges) that have a MeSH cross-reference.

\section{Limitations}

\begin{itemize}
  \item \textbf{Filtered edge set.} SemMedDB-KGX excludes low-support and
        constraint-violating predications, so recall describes this curated
        release rather than raw SemMedDB, and is likely lower than raw
        SemMedDB's.
  \item \textbf{Identifier mapping.} CUIs without a MeSH cross-reference, and
        CTD diseases identified only by OMIM, cannot be matched. Matching is
        on exact MeSH identifiers, with no credit for parent or child
        concepts in the MeSH hierarchy.
  \item \textbf{Direction and category.} Only chemical $\rightarrow$ disease
        edges are considered, and only for nodes whose \emph{most specific}
        category is in the configured sets. A chemical whose most specific
        category is outside these sets (for example, a protein-like category)
        is not considered.
  \item \textbf{Predicate mapping.} Predicate recall depends on the choice of
        $P(e)$ in Table~\ref{tab:predmap}. Entity recall gives an upper
        bound that does not depend on that choice.
\end{itemize}

\section{Implementation and reproducibility}

The analyses are implemented in a Snakemake (v9.27.0) workflow. The
predicate-frequency and recall queries run in DuckDB 1.5.5 (Python 3.12),
inside a pinned container image; the downloads run in a separate Debian-based
image containing \code{curl} and \code{tar}. JSON Lines were read with an
explicit column projection, which avoids loading the large embedded
supporting-text field, and with a 64\,MiB maximum record size, because
individual heavily cited edges exceed DuckDB's 16\,MiB default. DuckDB was
given an explicit memory limit and an on-disk spill directory so that it could
process the inputs out of core. The workflow ran on the UCSF CoreHPC Slurm
cluster. Because the compute nodes have no outbound network access, both
source datasets were downloaded on a login node beforehand. The workflow
writes the following outputs:
\code{predicate\_frequencies.csv} and
\code{predicate\_frequency\_lookup.tsv} for the frequency analysis, and
\code{ctd\_recall\_pairs.csv} (one row per CTD assertion in $C_M$, with its
match flags and the matched SemMedDB predicates) and
\code{ctd\_recall\_summary.csv} (recall by evidence type and pooled) for the
recall analysis.

\begin{thebibliography}{9}

\bibitem{kilicoglu2012}
H.~Kilicoglu, D.~Shin, M.~Fiszman, G.~Rosemblat, T.~C.~Rindflesch.
SemMedDB: a PubMed-scale repository of biomedical semantic predications.
\emph{Bioinformatics}, 28(23):3158--3160, 2012.

\bibitem{semmedkgx}
Translator-CATRAX. SemMedDB-KGX.
\url{https://github.com/Translator-CATRAX/SemMedDB-KGX}.

\bibitem{ctd}
A.~P.~Davis \emph{et al.}
Comparative Toxicogenomics Database (CTD).
\emph{Nucleic Acids Research}. \url{https://ctdbase.org}.

\bibitem{biolink}
D.~R.~Unni \emph{et al.}
Biolink Model: a universal schema for knowledge graphs in clinical,
biomedical, and translational science.
\emph{Clinical and Translational Science}, 15(8):1848--1855, 2022.

\bibitem{duckdb}
M.~Raasveldt, H.~M\"uhleisen.
DuckDB: an embeddable analytical database.
In \emph{Proceedings of SIGMOD}, 2019.

\bibitem{snakemake}
F.~M\"older \emph{et al.}
Sustainable data analysis with Snakemake.
\emph{F1000Research}, 10:33, 2021.

\end{thebibliography}

\end{document}
